The apparatus of Sections \ref{sec.pointProcesses} and \ref{sec.priceFormation} is developed
on flow that we generate, and a result obtained on generated flow is a result about the
generator.
This section turns to recorded data.

Four things are absent from the generated setting by construction, and each of them bounds
what can be asked of it.

The intensities do not read the book.
There is no queue-depletion feedback, so the arrival rates are indifferent to whether a queue
is about to empty; only the marks see the book, and that is the mechanism behind its
restoring force.
There is in particular no informed trading and therefore no \emph{adverse selection}.
In a traded market part of the flow is informed, and a resting limit order is then short an
option: it is filled when someone wishes to trade against it, which is disproportionately
when trading against it is right, and the defence available to whoever posted it is to
withdraw before that happens.
A queue about to be picked off is therefore a queue whose own side pulls away from it, and
depletion begets further depletion.
That is what makes $\queueImbalance[1]$ informative in a traded book:
a thin bid is thin because those posting on that side inferred something.
A kernel built for replenishment, as in Section \ref{sec.orderFlowModel}, does the reverse,
and the two cannot both be strong in one specification.

The marks are independent of the flow.
Definition \ref{def.orderFlowProcess} asks the sizes to be independent of the arrival times
and of the types, and Proposition \ref{prop.shortHorizonForecast} uses that.
A generator whose withdrawals draw against the size actually resting, or which drops a
withdrawal on an empty side, emits a thinned process rather than the Hawkes process it
started from, and the hypothesis fails first on the withdrawal types.

There is a single excitation timescale, the kernel being one exponential;
there is no intraday non-stationarity, so a window tuned in one regime is not tested against
another;
and there is no anchor for the price level, so the level wanders and the mean price change
over a session is a property of the seed.

\begin{remark}\label{remark.asymmetricReading}
 Any comparison of $\queueImbalance[1]$ with $\OFI$ conducted in the generated setting is
 therefore read asymmetrically.
 Whatever forecasting power $\queueImbalance[1]$ retains there is mechanical ---
 a depleted queue is one market order from clearing, and clearing moves the mid ---
 and none of it is informational.
 A result favourable to $\OFI$ is not evidence that $\OFI$ dominates $\queueImbalance[1]$ in
 a traded market; a result favourable to $\queueImbalance[1]$ is the stronger finding.
\end{remark}

On recorded data we estimate the kernel,
apply the diagnostics of Section \ref{sec.simulation} to recorded rather than simulated paths,
and ask the questions the generated setting cannot answer:
whether $\signedExcitation$ of \eqref{eq.signedExcitation} carries either sign in a traded
market, whether a single excitation timescale suffices to describe one,
and whether the sign of \eqref{eq.priceForecast} can be read from $\OFI$ alone.
